\chapter{Abdominal Surgery for Crohn's Disease}
\index{Abdominal Surgery for Crohn's Disease}

\section*{Definition and Overview}
Crohn's disease is a chronic, relapsing inflammatory condition of the digestive tract that can affect any region from the mouth to the anus, most commonly the terminal ileum and the beginning of the large bowel. Although medical therapy is the mainstay of management, surgery plays an essential role for a large proportion of patients. Abdominal surgery is used when the disease produces complications such as strictures, fistulae, or abscesses, or when medical treatment fails to control symptoms. Surgery cannot cure Crohn's disease, and the underlying inflammatory tendency remains, but it can effectively relieve symptoms, resolve complications, and restore quality of life. The guiding principle of operative management is to remove or repair the diseased segment while preserving as much healthy bowel as possible.

\section*{Epidemiology}
Crohn's disease is most commonly diagnosed in young adults between 15 and 35 years of age, with a second smaller peak in later life, and affects men and women fairly equally. Approximately 1 in 250 people in many countries live with inflammatory bowel disease, of which Crohn's disease is one of the two principal forms. Historically, between half and two-thirds of people with Crohn's disease required surgery at some point in their lives. The introduction of more effective medical therapies, including biologic agents, has reduced, but not eliminated, the need for surgery; it remains an important component of care for many patients, often within the first few years of diagnosis.

\section*{Aetiology and Pathophysiology}
The cause of Crohn's disease is incompletely understood but is thought to involve an abnormal immune response to gut flora in genetically susceptible individuals, with smoking as a major environmental trigger. Unlike the superficial inflammation of ulcerative colitis, the inflammation in Crohn's disease is transmural, affecting all layers of the bowel wall. This full-thickness involvement explains the distinctive complications that bring patients to surgery.

\begin{itemize}
    \item \textbf{Strictures:} fibrotic narrowing of the bowel causing cramps, bloating, and obstructive symptoms
    \item \textbf{Fistulae:} abnormal tracts between the bowel and adjacent bowel, bladder, vagina, or skin
    \item \textbf{Abscesses:} intra-abdominal collections of pus that fail to resolve with antibiotics or drainage
    \item \textbf{Perforation or severe haemorrhage} of the bowel
    \item \textbf{Failure of medical therapy} or unacceptable side effects of medications
    \item \textbf{Growth failure:} in children whose disease is difficult to control
    \item \textbf{Cancer prevention:} prophylactic surgery for long-standing, extensive colonic disease carrying an increased risk of colorectal cancer
\end{itemize}

\section*{Clinical Features}
The clinical picture after surgery includes the normal changes of recovery, which must be distinguished from complications.

\begin{itemize}
    \item Pain at the surgical incisions, which typically improves over the first few days
    \item Reduced appetite and a gradual return of bowel activity over several days
    \item A small amount of blood-stained fluid draining from the wound
    \item Significant tiredness and reduced stamina for several weeks
    \item Altered bowel habit, including looser or more frequent stools, which usually improves with time
    \item For people with a stoma, the need to learn stoma care and adapt to the appliance
\end{itemize}

\section*{Differential Diagnosis}
After surgery, the distinction between normal recovery and complications is critical.

\begin{itemize}
    \item Anastomotic leak versus a slow return of bowel function
    \item Intra-abdominal abscess versus normal postoperative swelling and discomfort
    \item Early disease recurrence versus adhesive obstruction
    \item Wound infection versus the normal inflammatory response to surgery
    \item Venous thromboembolism versus postoperative calf pain or breathlessness
    \item Urinary retention versus urinary tract infection
\end{itemize}

\section*{When to Seek Medical Care}
After surgery, the surgical team should be contacted urgently if any of the following develop, as they may indicate a serious complication.

\begin{itemize}
    \item Severe or worsening abdominal pain, or pain not controlled by prescribed analgesia
    \item Fever that does not settle
    \item Persistent vomiting, or inability to pass gas or open the bowels
    \item Wound redness, warmth, swelling, or discharge that is worsening
    \item Heavy or persistent bleeding from the wound or a drain
    \item Chest pain, shortness of breath, or pain and swelling in one calf
    \item Increasing abdominal swelling or signs of dehydration
\end{itemize}

\textbf{Contact your local emergency services for:} collapse, uncontrolled bleeding, sudden severe abdominal or chest pain, or difficulty breathing.

\section*{Diagnosis and Investigations}
Comprehensive assessment is required before surgery and again afterwards to plan follow-up.

\begin{itemize}
    \item \textbf{Pre-operative assessment:} detailed history, examination, medication review, nutritional assessment, and evaluation of steroid use
    \item \textbf{Colonoscopy:} to define the extent and activity of disease and to check for dysplasia
    \item \textbf{Magnetic resonance imaging (MRI) or CT:} to map strictures, fistulae, and abscesses for operative planning
    \item \textbf{Blood tests:} full blood count, inflammatory markers, nutritional indices, and blood group cross-matching
    \item \textbf{Anaesthetic assessment} and informed consent
    \item \textbf{Postoperative monitoring:} clinical review, blood tests, and later scheduled endoscopy or imaging to detect early endoscopic recurrence
\end{itemize}

\section*{Management}
The choice of procedure depends on the location and nature of the disease and on the patient's overall fitness.

\begin{itemize}
    \item \textbf{Pre-operative optimisation:} correction of nutrition, drainage of abscesses, and tapering of steroids before surgery to reduce complications
    \item \textbf{Strictureplasty:} widening of a narrowed segment without removing it, used to preserve bowel length
    \item \textbf{Bowel resection:} removal of the diseased segment with rejoining of healthy ends (anastomosis)
    \item \textbf{Colectomy:} removal of part or all of the colon, sometimes with the formation of an ileostomy or colostomy
    \item \textbf{Drainage of abscesses and definitive fistula surgery} where appropriate
    \item \textbf{Laparoscopic surgery:} increasingly preferred where feasible, offering a faster recovery and fewer wound complications
    \item \textbf{Postoperative medical therapy:} immunosuppressants or biologic agents commenced after surgery to reduce the risk of recurrence
    \item \textbf{Stoma care:} education and ongoing support from specialist stomal therapy nurses
    \item \textbf{Long-term follow-up:} regular review with scheduled endoscopy to detect and treat recurrence early
\end{itemize}

\section*{Complications}
\begin{itemize}
    \item Anastomotic leak leading to peritonitis or intra-abdominal abscess
    \item Wound infection and late incisional hernia
    \item Haemorrhage requiring transfusion or further surgery
    \item Venous thromboembolism involving the legs or lungs
    \item Bowel obstruction caused by postoperative adhesions
    \item Recurrence of Crohn's disease, most commonly at the anastomosis
    \item Short bowel syndrome, a rare but serious consequence of extensive or repeated resection
    \item Stoma-related problems, including skin irritation, prolapse, and leakage
\end{itemize}

\section*{Prognosis and Outcomes}
Surgery is highly effective at relieving the symptoms and complications that prompted it, and most people return to good health and full activity within weeks to months. However, Crohn's disease is not cured by resection, and the tendency to recur persists: endoscopic evidence of recurrence is present in a substantial proportion of patients within the first year, and many will eventually require further medical treatment or further surgery. Smoking cessation and the regular use of prescribed postoperative medications substantially reduce the risk of recurrence. With coordinated medical and surgical care, most people with Crohn's disease live full and productive lives.

\section*{Prevention and Self-Care}
\begin{itemize}
    \item Do not smoke, as smoking strongly increases the risk of recurrence after surgery
    \item Take prescribed medicines exactly as directed, since they reduce the chance of the disease returning
    \item Attend all follow-up appointments, including planned endoscopy
    \item Eat a healthy, balanced diet and obtain dietetic advice if weight or nutrition is a concern
    \item Resume exercise gradually, following the surgical team's advice on activity and lifting
    \item Report new symptoms such as abdominal pain, fever, or a change in bowel habit promptly
    \item If a stoma is present, make use of stomal therapy support and allow time to adapt
\end{itemize}

\section*{Sources and Further Reading}
The following organisations provide reliable information about Crohn's disease and the role of surgery in its management.

\begin{itemize}
    \item NHS. \textit{Crohn's disease: overview and treatment.} https://www.nhs.uk/conditions/crohns-disease/
    \item Mayo Clinic. \textit{Crohn's disease: symptoms, diagnosis, and treatment.} https://www.mayoclinic.org/diseases-conditions/crohns-disease/
    \item NICE. \textit{Crohn's disease: management guidance.} https://www.nice.org.uk/guidance/ng129
    \item Crohn's \& Colitis Foundation. \textit{Surgery for Crohn's disease.} https://www.crohnscolitisfoundation.org/
    \item MSD Manual. \textit{Crohn disease: professional reference.} https://www.msdmanuals.com/
    \item MedlinePlus. \textit{Crohn disease health information.} https://medlineplus.gov/crohnsdisease.html
\end{itemize}
